\section{Functions}

\begin{definition}
    A \textbf{function} $f: D \to C$ is a map from a domian $D$ to a codomain $C$ s.t. for each $x \in D$ there is exactly one $y \in C$ st. $f(x) = y$ (well defined).
\end{definition}

\begin{definition}
    A function $f$ is called \textbf{bouned} if there exists $M \in \mathbb{R}$ such that $|f(x)| \leq(\geq) M$ for all $x \in D$.
\end{definition}

\begin{definition}
    A function $f$ is called \textbf{injective} iff 
    $$\forall x_1, x_2 \in A : f(x_1) = f(x_2) \implies x_1 = x_2$$
    \textbf{surjective} iff 
    $$\forall y \in B \exists x \in A : f(x) = y$$
    and bijective iff $f$ is injective and surjective.
\end{definition}

\begin{definition}
    The \textbf{restriction} of $f$ under $D' \subset \mathbb{D}$ := 
    $$f\upharpoonleft_{D'}: D' \to C, x \mapsto f(x) \quad \forall x \in \mathbb{D}$$
\end{definition}

\begin{definition}
    Let $f: D \to C$ and $g: C \to B$ be functions then
    $$g \circ f: D \to B \quad \text{with}$$
    $$g \circ f(x) = g(f(x)) \quad \forall x \in D$$
    is called the \textbf{composition} of $f$ and $g$.
\end{definition}

\begin{definition}
    The \textbf{identity function} $id_D: D \to D$ := $id_D(x) = x \quad \forall x \in D$.
\end{definition}

\begin{definition}
    If a function $f: D \to C$ is bijective then there exists a unique function $f^{-1}: C \to D$ such that
    $$ f(f^{-1}(y)) = id_C = y \quad \forall y \in C$$
    and
    $$ f^{-1}(f(x)) = id_D = x \quad \forall x \in D$$
    $f^{-1}$ is called the \textbf{inverse function} of $f$.
\end{definition}

\begin{definition}
    A function $f: \mathbb{R} \to \mathbb{R}$ is called
    \textbf{even} iff 
    $$\forall x \in \mathbb{R} : f(-x) = f(x)$$
    and \textbf{odd} iff 
    $$\forall x \in \mathbb{R} : f(-x) = -f(x)$$
\end{definition}

\begin{definition}
    A function  $f: \mathbb{R} \to \mathbb{R}$ is called \textbf{(strictly) monotonically increasing} iff 
    $$\forall x_1, x_2 \in \mathbb{R} : x_1 < x_2 \implies f(x_1) < f(x_2)$$
    and \textbf{(strictly) monotonically decreasing} iff 
    $$\forall x_1, x_2 \in \mathbb{R} : x_1 < x_2 \implies f(x_1) > f(x_2)$$
\end{definition}

\begin{theorem}
    Every strictly monotonic function is injective.
\end{theorem}


\subsection{Limits in the context of functions}

\begin{definition}
    Let $f: \mathbb{D} \to \mathbb{R}$ be a function satisfying $\mathbb{D} \cap (x_0 - \delta, x_0 + \delta) \neq \emptyset$ for all $\delta > 0$ and $x_0 \in \mathbb{R}$. We say that $\lim_{x \to x_0} f(x) = L$ iff 
    $$\forall \varepsilon > 0 \exists \delta > 0  \text{s.t.}$$
    $$x \in \mathbb{D} \cap (x_0 - \delta, x_0 + \delta) \implies |f(x) - L| < \varepsilon$$
    The limit $L$ is unique.
\end{definition}

\begin{definition}
    A function $f: \mathbb{D} \to \mathbb{R}$ has in $x_0$ a
    \textbf{right-sided limit} $\quad \lim_{x \searrow x_0} f(x) = \lim_{x \to x_0^+} f(x) = L$ iff 
    $$\forall \varepsilon > 0 \exists \delta > 0  \quad \text{s.t.}$$
    $$\forall x \in \mathbb{D} (x_0 < x < x_0 + \delta \implies |f(x) - L| < \varepsilon)$$
\end{definition}

\begin{definition}
    A function $f: \mathbb{D} \to \mathbb{R}$ has a \textbf{limit for x approaching infinity} $\lim_{x \to \infty} f(x) = L$ iff 
    $$\forall \varepsilon > 0 \exists M > 0 \quad \text{s.t.}$$
    $$\forall x \in \mathbb{D} (x > M \implies |f(x) - L| < \varepsilon)$$
\end{definition}

\begin{definition}
    A function $f: \mathbb{D} \to \mathbb{R}$ has in $x_0$ a \textbf{improper limit approaching infinity} $\lim_{x \to x_0} f(x) = \infty$ iff 
    $$\forall N > 0 \exists \delta > 0 \quad \text{s.t.}$$
    $$forall x \in \mathbb{D} (0 < |x - x_0| < \delta \implies f(x) > N)$$
\end{definition}

\begin{theorem}
    Let $\lim_{x \to x_0} f(x) = K (\neq \pm \infty)$ and $\lim_{x \to x_0} g(x) = L (\neq \pm \infty)$
    $$\lim_{x \to x_0} (f(x) + g(x)) = K + L$$
    $$\lim_{x \to x_0} (f(x) - g(x)) = K - L$$
    $$\lim_{x \to x_0} (f(x) g(x)) = K L$$
    $$\lim_{x \to x_0} (f(x) / g(x)) = K / L$$
    $$\lim_{x \to x_0} (A f(x)) = A K$$
    $$\lim_{x \to x_0} (f(x)/g(x)) = K/L \text{if} L \neq 0$$
\end{theorem}

\subsection{Continuity of functions}

\begin{definition}
    A function $f: D \to \mathbb{R}$ is called \textbf{continous} at $x_0$ if
    $$\forall \epsilon \gt 0 \exists \delta \gt 0 \text{ such that }$$
    $$\forall x \in I (\mid x - x_0 \mid \lt \delta \implies \mid f(x) - f(c) \mid \gt \epsilon)$$
    The function is called continuous if it is continuous at every point in its domain.
\end{definition}

\begin{remark}
    The \textbf{intuition} behind this definition is that for any given distance $\epsilon > 0$ around $f(x_0)$, we can find a $\delta > 0$ such that for all $x \in D$ that are at least $\delta$ close to $x_0$, we have $|f(x) - f(x_0)| < \epsilon$. 
    So for any bound on the output ($\epsilon$) there exists a bound on the input ($\delta$) such that all elements in the input bound $(x_0 - \delta, x_0 + \delta)$ are mapped to the output bound $(f(x_0) - \epsilon, f(x_0) + \epsilon)$.
\end{remark}

\begin{definition}
    Let $f: D \to \mathbb{R}$ be a function. Then $f$ is called \textbf{left(right)-continuous} at $x_0$ if
    $$\lim_{x \to x_0^{-(+)}} f(x) = f(x_0)$$
\end{definition}

\begin{theorem}
    Let $f: D \subset \mathbb{D}(f) \subset \mathbb{R}$ be a function and $x_0 \in D$. $f$ is continuous at $x_0$ iff for every sequence $(x_n)_{n\in\mathbb{N}} \subset D$ with $\lim_{n\to\infty} x_n = x_0$ it holds that $\lim_{n\to\infty} f(x_n) = f(x_0)$.
\end{theorem}

\begin{remark}
    Continuity can also be viewed in terms of intervals, where a function $f$ is continuous iff for every open interval $I \subset \mathbb{R}$ it holds that $f^{-1}(I)$ is open.
\end{remark}

\begin{lemma}
    Let $f, g$ be continuous functions then
    $$f+g, \quad f-g, \quad f \cdot g, \quad c \cdot f, \quad \frac{f}{g}, \quad f \circ g$$
    are continuous
\end{lemma}

\begin{theorem}
    Let $I$ be an interval and let $f: I \to J$ be a continuous and bijective function. Then $J$ is an interval too and the inverse function $f^{-1}: J \to I$ is continuous.
\end{theorem}

\begin{theorem}
    Let $f: D \to \mathbb{R}$ be a continuous function. Let $c \in [f(a), f(b)]$ then there exists a $x_0 \in [a, b]$ such that $f(x_0) = c$.
\end{theorem}

\begin{lemma}
    Let $[a, b]$ be a compact interval and let $(x_n)_{n\in\mathbb{N}} \subset [a, b]$ be a sequence. Then there exists a subsequence $(x_{n_k})_{k\in\mathbb{N}}$ such that 
    $$ \lim_{k\to\infty} x_{n_k} = x_0 \quad \text{ for some } x_0 \in [a, b]. $$
\end{lemma}

\begin{definition}
    Let $f: D \subset \mathbb{R} \to \mathbb{R}$ be a function
    If for $x_0 \in D$ it holds that $f(x_0) \geq f(x) \forall x \in D$, then $f$ has a local maximum at $x_0$.
    If for $x_0 \in D$ it holds that $f(x_0) \leq f(x) \forall x \in D$, then $f$ has a local minimum at $x_0$.
\end{definition}

\begin{theorem}
    Let $f: I \subset \mathbb{R} \to \mathbb{R}$ be a continous function and $I$ be a compact interval. Then $f$ is bounded and attains its maximum and minimum value on $I$. Hence
    $$\exists x_{min}, x_{max} \in I \text{ such that }$$
    $$ f(x_{min}) \leq f(x) \leq f(x_{max}) \quad \forall x \in I$$
\end{theorem}

\subsection{Trigonometric Identities}
\begin{center}
\scriptsize
\setlength{\tabcolsep}{2.5pt}
\begin{tabular}{|l|l|}
\hline
\textbf{Category} & \textbf{Formula} \\ \hline
Pythagoras & $\sin^2 x + \cos^2 x = 1$ \\
& $1 + \tan^2 x = \frac{1}{\cos^2 x}$ \\ \hline
Addition & $\sin(x \pm y) = \sin x \cos y \pm \cos x \sin y$ \\
& $\cos(x \pm y) = \cos x \cos y \mp \sin x \sin y$ \\
& $\tan(x \pm y) = \frac{\tan x \pm \tan y}{1 \mp \tan x \tan y}$ \\ \hline
Double Angle & $\sin(2x) = 2 \sin x \cos x$ \\
& $\cos(2x) = \cos^2 x - \sin^2 x$ \\
& $\phantom{\cos(2x)} = 2 \cos^2 x - 1 = 1 - 2 \sin^2 x$ \\ \hline
Half Angle & $\sin^2 x = \frac{1 - \cos(2x)}{2}$ \\
& $\cos^2 x = \frac{1 + \cos(2x)}{2}$ \\ \hline
Subst. trick & $\tan(x/2) = t \implies$ \\
& $\sin x = \frac{2t}{1+t^2}, \cos x = \frac{1-t^2}{1+t^2}$ \\
& $dx = \frac{2dt}{1+t^2}$ \\ \hline
Prod-to-Sum & $\sin x \sin y = \frac{1}{2}(\cos(x-y) - \cos(x+y))$ \\
& $\cos x \cos y = \frac{1}{2}(\cos(x-y) + \cos(x+y))$ \\
& $\sin x \cos y = \frac{1}{2}(\sin(x+y) + \sin(x-y))$ \\ \hline
\end{tabular}
\end{center}

\subsection{Hyperbolic Functions}
\begin{itemize}
    \item $\cosh x = \frac{e^x + e^{-x}}{2}, \quad \sinh x = \frac{e^x - e^{-x}}{2}$
    \item $\tanh x = \frac{\sinh x}{\cosh x} = \frac{e^x - e^{-x}}{e^x + e^{-x}}$
    \item $\cosh^2 x - \sinh^2 x = 1$
    \item $\sinh(2x) = 2 \sinh x \cosh x, \quad \cosh(2x) = \cosh^2 x + \sinh^2 x$
    \item $(\sinh x)' = \cosh x, \quad (\cosh x)' = \sinh x, \quad (\tanh x)' = \frac{1}{\cosh^2 x}$
    \item $\text{arsinh}(x) = \ln(x + \sqrt{x^2+1}), \quad \text{arcosh}(x) = \ln(x + \sqrt{x^2-1})$
\end{itemize}
